\chapter{Resultados}\label{cap:resultados}

Este capítulo presenta los resultados obtenidos: el conjunto de datos y su \textit{ground truth}, la evaluación de los SLM, la sensibilidad al tamaño de \textit{batch}, la fiabilidad de ejecución y el estado de los componentes de notificación, recomendación e inyección de errores. Todas las cifras provienen de los reportes del proyecto técnico (\texttt{analysis/batch\_size\_matrix/}, \texttt{pseudo\_labeling/}) o, cuando se indica, de un recálculo propio de solo lectura sobre esos archivos (\texttt{notes/eval\_final\_gt.py}). Donde un resultado no existe todavía, se indica con un comentario \texttt{TODO}.

\section{Conjunto de datos (objetivo específico 1)}\label{sec:res_dataset}

\subsection{Prefiltrado}

De 27{,}752{,}836 eventos del corpus \texttt{ms-users} se conservaron 223{,}104 (0.804\,\%), organizados en 4{,}463 \textit{batches} de 50 eventos (4{,}462 completos y uno de 4). % TODO: tabla con el desglose del filter_report.txt (eventos descartados por motivo) y distribución por nivel/logger de los eventos conservados.

\subsection{Preetiquetado por consenso}

La votación de tres LLM sobre cada \textit{batch} produjo 1{,}370 \textit{batches} error (30.70\,\%) y 3{,}093 normales por mayoría, con $\kappa$ de Fleiss de 0.886 (Sección~\ref{subsec:met_kappa}). Los votos individuales, el tiempo y el consumo de \textit{tokens} están en la Tabla~\ref{tab:met_votos}. El consenso coincide con la referencia derivada de las etiquetas por evento (error si algún evento lo es) en 4{,}400 de 4{,}463 \textit{batches} (98.59\,\%).

\subsection{Revisión experta}

La Tabla~\ref{tab:res_expertos} muestra el resultado de la segunda revisión de la muestra aleatoria estratificada de 220 \textit{batches}, hecha con el criterio acordado (Tabla~\ref{tab:met_criterio}). Del total, 63 etiquetas pasaron de error a normal respecto de la primera revisión y ninguna de normal a error.

\begin{table}[H]
\centering
\caption{Segunda revisión experta de la muestra aleatoria estratificada. IC: intervalo de Wilson al 95\,\%.}
\label{tab:res_expertos}
\small
\begin{tabular}{@{}lrrrr@{}}
\toprule
\textbf{Estrato (según los LLM)} & \textbf{\textit{Batches} revisados} & \textbf{Error según experto} & \textbf{IC 95\,\%} & \textbf{Contradice al consenso} \\
\midrule
Error unánime (3/3) & 88 & 11 (12.5\,\%) & 7.1--21.0\,\% & 77/88 (87.5\,\%) \\
Desacuerdo (1 o 2 votos) & 88 & 6 (6.8\,\%) & 3.2--14.1\,\% & 43/88 (48.9\,\%) \\
Normal unánime (0/3) & 44 & 0 (0\,\%) & 0--8.0\,\% & 0/44 \\
\bottomrule
\end{tabular}
\normalsize
\end{table}

El hallazgo central es que \textbf{el consenso de tres LLM sobreestima el error}: solo el 12.5\,\% de los \textit{batches} en los que los tres modelos coinciden en ``error'' lo son según los expertos. Los modelos marcaban como error condiciones que los revisores consideran funcionales (política de contraseñas, respuestas 404/406, usuario no encontrado), y en parte esto se explica porque el \textit{prompt} de la votación las incluye como errores (Apéndice~\ref{ap:prompt_votacion}). En total los expertos revisaron 329 \textit{batches} (220 de la muestra y 109 de rondas dirigidas): 39 error y 290 normales.

\subsection{\textit{Ground truth} final}

La propagación del criterio produjo etiquetas para los 4{,}463 \textit{batches}, sin ninguno pendiente (Tabla~\ref{tab:res_origen}). El \textit{ground truth} final contiene \textbf{130 \textit{batches} error (2.91\,\%)}: 39 por revisión directa y 91 por reglas; la similitud solo asigna normal. Por estrato original, hay 120 errores en error unánime (de 1{,}203), 10 en desacuerdo (de 324) y 0 en normal unánime. El consenso preliminar y el \textit{ground truth} final coinciden en 3{,}221 de 4{,}463 \textit{batches} (72.2\,\%).

\begin{table}[H]
\centering
\caption{Origen de la etiqueta en el \textit{ground truth} final.}
\label{tab:res_origen}
\small
\begin{tabular}{@{}lrrr@{}}
\toprule
\textbf{Origen (\texttt{label\_source})} & \textbf{\textit{Batches}} & \textbf{Error} & \textbf{Normal} \\
\midrule
Ningún LLM votó error & 2{,}892 & 0 & 2{,}892 \\
Similitud con \textit{batch} revisado & 1{,}093 & 0 & 1{,}093 \\
Experto & 329 & 39 & 290 \\
Regla ``AS400 normal'' & 58 & 0 & 58 \\
Regla de error & 55 & 55 & 0 \\
Regla de error + similitud & 31 & 31 & 0 \\
Regla de error sobre similitud & 5 & 5 & 0 \\
\textbf{Total} & \textbf{4{,}463} & \textbf{130} & \textbf{4{,}333} \\
\bottomrule
\end{tabular}
\normalsize
\end{table}

La estimación por muestreo aleatorio extrapola unos 172 errores (3.9\,\%) con un intervalo muy amplio (96--534), de modo que el valor final de 130 cae dentro del intervalo pero por debajo de la estimación puntual: el \textit{ground truth} podría subestimar algunos errores. La validación \textit{leave-one-out} de las reglas sobre los 285 \textit{batches} revisados con voto de error da TP $=17$, FP $=0$, FN $=8$, TN $=174$ y 86 sin resolver (precisión 1.00, \textit{recall} 0.44), pero es optimista porque los patrones salen de las mismas etiquetas.
% TODO: tabla de tipos de error del GT final (por regla/tipo de evento) para cerrar el objetivo 1 "clasificados por tipos".
% TODO: acuerdo entre expertos (no se midió) y muestra independiente de comprobación de lo propagado: están pendientes (GROUND_TRUTH_BATCHES.md, sección 10).

\section{Evaluación de los SLM (objetivo específico 3)}\label{sec:res_slm}

\subsection{Resultados contra la referencia por consenso, $N=50$}\label{subsec:res_n50}

La Tabla~\ref{tab:res_n50} presenta las métricas de los ocho SLM con $N=50$ frente a la colección de referencia \texttt{gt\_\_final\_\_consensus\_is\_error\_\_n50} (consenso por evento; 1{,}370 \textit{batches} positivos de 4{,}463, es decir 30.7\,\%), con $k=1$. En todas las corridas de $N=50$ el vecino más cercano es el mismo \textit{batch} (\texttt{same\_batch\_top1\_rate}$=1.0$), de modo que la comparación equivale a una comparación posicional. \textbf{Esta referencia no es el \textit{ground truth} corregido por expertos}; ver la Sección~\ref{subsec:res_gtfinal}.

\begin{table}[H]
\centering
\caption{Métricas de los SLM con $N=50$ contra la referencia por consenso (kNN, $k=1$). \textit{Batches}: número evaluado (válidos).}
\label{tab:res_n50}
\small
\begin{tabular}{@{}lrrrrrr@{}}
\toprule
\textbf{Modelo} & \textbf{\textit{Batches}} & \textbf{Acc.} & \textbf{Prec.} & \textbf{Recall} & \textbf{$F_1$} & \textbf{AUC} \\
\midrule
granite-4.1-3b & 4{,}463 & 0.971 & 0.990 & 0.915 & 0.951 & 0.997 \\
Phi-3.5-mini-instruct & 4{,}461 & 0.949 & 0.994 & 0.840 & 0.911 & 0.997 \\
gemma-3-4b-it & 4{,}455 & 0.682 & 0.491 & 0.956 & 0.649 & 0.948 \\
Qwen3.5-9B (Q6\_K) & 4{,}463 & 0.828 & 0.919 & 0.482 & 0.632 & 0.978 \\
Ling-3.0-tiny (Q4\_K\_M) & 4{,}460 & 0.526 & 0.382 & 0.880 & 0.533 & 0.909 \\
Ministral-3-3B (BF16) & 4{,}312 & 0.424 & 0.341 & 0.983 & 0.506 & 0.872 \\
llama-3.2-1B-log-analyzer (Q4\_K\_M) & 4{,}441 & 0.690 & 0.432 & 0.047 & 0.085 & 0.573 \\
SmolLM2-1.7B (parcial) & 2{,}200 & 0.615 & 0.000 & 0.000 & 0.000 & 0.470 \\
\bottomrule
\end{tabular}
\normalsize
\end{table}

% TODO: convertir y copiar a figs/resultados/ las curvas ROC y matrices de confusión (analysis/batch_size_matrix/comparison/*/roc_knn.png, confusion_knn.png) y las figuras SVG f1_por_modelo_dotplot.svg y f1_vs_recall_dumbbell(2).svg (raíz del proyecto técnico) a PDF; insertar con \includegraphics.

Contra esta referencia, granite-4.1-3b y Phi-3.5-mini son los mejores (\textit{accuracy} 0.971 y 0.949, $F_1$ 0.951 y 0.911, AUC 0.997 ambos). Gemma, Ling y Ministral sobredetectan error: \textit{recall} alto (0.956, 0.880, 0.983) y precisión baja (0.491, 0.382, 0.341), lo que en operación equivale a generar fatiga de alertas. El modelo llama-3.2-1B-log-analyzer casi nunca marca error (\textit{recall} 0.047). Qwen3.5-9B muestra el comportamiento opuesto a gemma: precisión alta (0.919) y \textit{recall} bajo (0.482), aunque con AUC de 0.978, lo que sugiere que un umbral distinto sobre $P(\text{error})$ mejoraría su \textit{recall}. SmolLM2-1.7B solo tiene 2{,}200 \textit{batches} evaluables porque 2{,}263 fallaron, y no predijo ningún error en ellos.

\subsection{Recálculo contra el \textit{ground truth} final}\label{subsec:res_gtfinal}

Como la referencia de la sección anterior sobreestima el error, se recalcularon las métricas de las corridas de $N=50$ contra la etiqueta final (\texttt{final\_is\_error}, 130 positivos) alineando por posición del \textit{batch}; los \textit{batches} inválidos se excluyen. Los conteos de \textit{batches} predichos como error coinciden con los de los reportes de la matriz, lo que confirma la alineación. El resultado se muestra en la Tabla~\ref{tab:res_n50_final}. Este cálculo es propio (no proviene de \texttt{collection\_comparison.py}).

\begin{table}[H]
\centering
\caption{Métricas con $N=50$ contra el \textit{ground truth} final por \textit{batch} (130 errores en 4{,}463). TP, FP, TN, FN con \textit{error} como clase positiva. Recálculo propio.}
\label{tab:res_n50_final}
\small
\begin{tabular}{@{}lrrrrrrrrr@{}}
\toprule
\textbf{Modelo} & \textbf{Válidos} & \textbf{Acc.} & \textbf{Prec.} & \textbf{Recall} & \textbf{$F_1$} & \textbf{AUC} & \textbf{TP} & \textbf{FP} & \textbf{FN} \\
\midrule
Qwen3.5-9B & 4{,}463 & 0.866 & 0.174 & 0.962 & 0.295 & 0.987 & 125 & 593 & 5 \\
granite-4.1-3b & 4{,}463 & 0.743 & 0.099 & 0.962 & 0.179 & 0.932 & 125 & 1{,}142 & 5 \\
Phi-3.5-mini & 4{,}461 & 0.762 & 0.097 & 0.868 & 0.174 & 0.904 & 112 & 1{,}045 & 17 \\
gemma-3-4b-it & 4{,}455 & 0.429 & 0.047 & 0.961 & 0.089 & 0.841 & 124 & 2{,}541 & 5 \\
llama-3.2-1B-log-analyzer & 4{,}441 & 0.943 & 0.081 & 0.094 & 0.087 & 0.549 & 12 & 136 & 116 \\
Ling-3.0-tiny & 4{,}460 & 0.319 & 0.040 & 0.962 & 0.076 & 0.933 & 125 & 3{,}031 & 5 \\
Ministral-3-3B & 4{,}312 & 0.161 & 0.032 & 0.984 & 0.062 & 0.876 & 120 & 3{,}616 & 2 \\
SmolLM2-1.7B (parcial) & 2{,}200 & 0.957 & 0.000 & 0.000 & 0.000 & 0.361 & 0 & 0 & 95 \\
\bottomrule
\end{tabular}
\normalsize
\end{table}

El cambio de referencia modifica sustancialmente la lectura. Casi todos los modelos detectan más del 86\,\% de los 130 errores reales (\textit{recall} de 0.868 a 0.984), pero su precisión es muy baja (0.032 a 0.174): señalan como error muchos \textit{batches} que los expertos consideran normales. Para granite, por ejemplo, de 1{,}267 \textit{batches} señalados solo 125 son errores según el \textit{ground truth} final, de modo que su precisión no puede superar 0.103. El mejor $F_1$ pasa a ser el de Qwen3.5-9B (0.295), que es también el de mayor AUC (0.987). SmolLM2 y llama-3.2-1B-log-analyzer logran \textit{accuracy} alta (0.957 y 0.943) solo porque casi nunca predicen error, lo que ilustra por qué \textit{accuracy} es insuficiente con 2.91\,\% de positivos. Los SLM de mejor rendimiento contra el consenso de LLM (granite, Phi) reproducen el criterio de los LLM de votación, que los expertos corrigieron; esto sugiere que la elevada concordancia con el consenso no implica concordancia con el criterio del dominio, en línea con el hallazgo de la sección anterior.
% TODO: verificar este recálculo con collection_comparison.py contra una colección construida desde ground_truth_batches_final.jsonl (build_collections.py --ground-truth ...), que es el procedimiento oficial; las cifras deberían coincidir en N=50 por ser top-1 = mismo batch.
% TODO: añadir curvas ROC y matrices de confusión contra el GT final; analizar los 1,142 falsos positivos de granite por tipo de evento (¿4xx, contraseñas, usuario no encontrado?) para verificar si el prompt no logra transmitir el criterio.

\subsection{Sensibilidad al tamaño de \textit{batch}}\label{subsec:res_tamano}

Los modelos se ejecutaron también con $N=30$ y $N=20$ y se compararon contra la referencia de 50 (mismo consenso por evento). La Tabla~\ref{tab:res_tamano} reporta el $F_1$ del vecino más cercano y el de la métrica por solapamiento de índices (\texttt{index\_overlap}).

\begin{table}[H]
\centering
\caption{$F_1$ por tamaño de \textit{batch} contra la referencia de consenso de 50 (primera cifra: kNN $k=1$; entre paréntesis: solapamiento de índices). ``--'': corrida pendiente.}
\label{tab:res_tamano}
\small
\begin{tabular}{@{}lccc@{}}
\toprule
\textbf{Modelo} & \textbf{$N=50$} & \textbf{$N=30$} & \textbf{$N=20$} \\
\midrule
granite-4.1-3b & 0.951 (0.951) & 0.553 (0.695) & 0.487 (0.584) \\
Phi-3.5-mini & 0.911 (0.911) & 0.415 (0.493) & 0.402 (0.480) \\
Qwen3.5-9B & 0.632 (0.632) & 0.347 (0.434) & -- \\
gemma-3-4b-it & 0.649 (0.649) & 0.550 (0.663) & -- \\
Ling-3.0-tiny & 0.533 (0.533) & 0.449 (0.518) & 0.438 (0.486) \\
Ministral-3-3B & 0.506 (0.506) & 0.477 (0.554) & -- \\
llama-3.2-1B-log-analyzer & 0.085 (0.085) & 0.124 (0.140) & 0.298 (0.335) \\
\bottomrule
\end{tabular}
\normalsize
\end{table}

Para los dos mejores modelos el $F_1$ cae de forma marcada al reducir el \textit{batch} (granite: 0.951, 0.553, 0.487), por lo que no hay evidencia de que \textit{batches} menores a 50 mejoren la detección. La interpretación debe ser cautelosa: al comparar \textit{batches} de tamaño distinto con una referencia de 50, la alineación semántica cambia (la similitud media del vecino más cercano es 0.994 y la mínima 0.441) y la proporción de positivos de la referencia efectiva cambia (por ejemplo, 2{,}242 positivos en $N=30$ y 3{,}338 en $N=20$ para granite), de modo que la caída refleja en parte el método de comparación y no solo el modelo. % TODO: SmolLM2 N=30/N=20 y gemma/Qwen/Ministral N=20 siguen pendientes (RESUMEN_CONTEXTO_TESIS.md, sección 5); completar la tabla. Recalcular con el GT final.

\subsection{Fiabilidad de ejecución y tiempo}\label{subsec:res_fiabilidad}

Con $N=50$ las corridas terminaron con 4{,}463 \textit{batches} procesados, pero algunos no produjeron respuesta válida: Ministral-3-3B 151 \textit{batches} inválidos (183 con $N=30$), llama-3.2-1B-log-analyzer 22, gemma-3-4b-it 8 fallidos, Ling-3.0-tiny 3 fallidos, Phi-3.5-mini 2 fallidos y SmolLM2-1.7B 2{,}263 fallidos (la mitad del corpus). La causa dominante en Ministral es un JSON casi válido (saltos de línea dentro de cadenas, comillas sin escapar); el \textit{parser} tolerante y \texttt{reparse\_run.py} atienden ese caso. % TODO: reparsear los inválidos de Ministral y llama con evaluation/reparse_run.py y reportar el efecto en las métricas; rehacer SmolLM2 N=50.

Los tiempos totales de las corridas de $N=50$ van de 6{,}442~s (granite-4.1-3b, 1.8 h) a 22{,}649~s (gemma-3-4b-it) y 51{,}188~s (SmolLM2-1.7B), con Qwen3.5-9B en 10{,}640~s. No son comparables directamente porque las corridas se hicieron en dos máquinas con GPU distinta (NVIDIA y AMD) y distinto motor (Ollama o Transformers). % TODO: tabla de tiempos por batch normalizados por máquina/backend y uso de VRAM.

\subsection{Experimentos previos con prompt y caché}\label{sec:res_previos}

Antes de la matriz actual se compararon variantes del \textit{prompt} y del caché semántico con \textit{evaluación posicional} contra el \textit{ground truth} por evento (consenso, agregado por \textit{batch} de 50; criterio anterior a la calibración experta). Con granite-4.1-3b: la variante base con caché dio $F_1$ de 0.549, el \textit{prompt} ajustado con caché 0.639 y el \textit{prompt} ajustado sin caché 0.897 (\textit{accuracy} 0.934, AUC 0.992). Con gemma-3-4b-it, el \textit{prompt} ajustado sin caché dio $F_1$ de 0.551 y con caché 0.521. Estas cifras sugieren que el \textit{prompt} aporta más que el caché y que el caché reduce la calidad al reutilizar clasificaciones; la matriz final se ejecutó sin caché (\texttt{embedding\_hits}$=0$). % TODO: confirmar el significado exacto de cada carpeta (base_and_embedding / prompt_no_embedding / prompt_and_embedding) con los autores y la configuración en run_<id>.json; las rutas del eval_report.json apuntan a otra máquina, y el prompt "base" no está documentado en el repositorio.

\section{Ajuste fino LoRA (objetivo específico 2)}\label{sec:res_lora}

El entrenamiento LoRA ($r=8$, $\alpha=32$, 3 épocas, 1{,}200 pasos) registró una pérdida de evaluación de 0.0071 al final de la primera época y de 0.0096 al final de la segunda, con pérdida de entrenamiento de 0.518 (paso 500) a 0.0037 (paso 1{,}000). Los \textit{checkpoints} existentes (500, 1{,}200) se entrenaron sobre un conjunto de 10{,}001 registros con solo 39 positivos, que no parece provenir de \texttt{ms-users} (Sección~\ref{sec:met_finetuning}). % TODO: no existe evaluación del adaptador ajustado contra el GT por batch; ejecutarla (y compararla con el modelo base) o declarar el ajuste fino como trabajo futuro. La pérdida de evaluación de 0.0096 con 39 positivos de 10,001 sugiere que el modelo aprendió la clase mayoritaria; verificar con métricas por clase.

\section{Notificación, recomendación e inyección de errores (objetivos 4 y 5)}\label{sec:res_reaccion}

\subsection{Implementación}

El \textit{pipeline} cuenta con notificación a Slack, recomendador local, inyección de errores desde el \textit{dashboard} y \textit{dashboard} de seguimiento, todos cubiertos por pruebas unitarias (Sección~\ref{sec:met_pruebas}). % TODO: capturas de pantalla del dashboard y de una alerta de Slack; ejemplo real de recomendación (recommendations.jsonl) para un batch de error del GT.

\subsection{Simulación de errores}

% TODO: no existen resultados de simulación. Las plantillas de inyección son placeholders y el protocolo no está definido. Una vez definidos, reportar por tipo de error inyectado: tasa de detección, latencia hasta la alerta y calidad de la recomendación; con al menos el modelo de mayor F1 contra el GT final (Qwen3.5-9B) y uno de los más rápidos (granite-4.1-3b).
Para el objetivo 5 no se dispone todavía de evidencia experimental; este resultado queda pendiente.
